% !PRINTED 66
\weberpage{66}
%<<CONT

\(1,2,\ldots,n\), in a definite way, \(n\) real numerical values \(a_1,a_2,\ldots,a_n\), and from these numerical values form the \(\frac{n(n-1)}2\) differences
\[
 a_1-a_2,\quad a_1-a_3,\ldots,
\]
where the lower index is to belong to the minuend. The product of differences
\begin{equation}
 P=(a_1-a_2)(a_1-a_3)\cdots(a_1-a_n)
    (a_2-a_3)\cdots(a_2-a_n)\cdots(a_{n-1}-a_n)
\tag{1}
\end{equation}
will, when \(a_1,a_2,\ldots,a_n\) are mutually distinct numerical values, have a value different from zero, for example a positive one if \(a_1>a_2>a_3>\cdots>a_n\) was assumed.

If now we interchange the indices \(1,2,3,\ldots,n\) in some way, and so pass, say, from \(\mathfrak A\) to \(\mathfrak A'\), then \(P\) passes into
\begin{equation}
\begin{aligned}
 P'={}&(a_{\alpha_1}-a_{\alpha_2})(a_{\alpha_1}-a_{\alpha_3})\cdots
       (a_{\alpha_1}-a_{\alpha_n})\\
     &\quad\cdot(a_{\alpha_2}-a_{\alpha_3})\cdots(a_{\alpha_2}-a_{\alpha_n})
       \cdots(a_{\alpha_{n-1}}-a_{\alpha_n}),
\end{aligned}
\tag{2}
\end{equation}
and, apart from the sign, this product consists of the same factors as \(P\). Thus \(P'\) is either equal to \(P\) or opposite to \(P\).

\gesperrt{I. We count the arrangement \(\mathfrak A'\), and therefore also the permutation which transforms \(\mathfrak A\) into \(\mathfrak A'\), as of the first or of the second kind according as \(P\) is equal or opposite to \(P'\); hence \(\mathfrak A\) itself belongs to the first kind.}

\gesperrt{II. By a simple transposition \((h,k)\), in which \(h,k\) denote any two of the digits \(1,2,\ldots,n\), both \(P\) and \(P'\) change their sign.}

For the factors which contain neither \(h\) nor \(k\) are not affected by this transposition; furthermore, in \(P\) and \(P'\) there is the factor \(\pm(a_h-a_k)\), and the pairs of factors
\[
 \pm(a_h-a_\nu)(a_k-a_\nu),
\]
where \(\nu\) runs through the series \(1,2,\ldots,n\) with the exception of \(h,k\). The first factor changes its sign, while the pair of factors remains unchanged under the transposition \((h,k)\). Hence follows:
%CONT>>
